\section{Conceptual foundations: multi-energy systems, surrogate models and optimization bottlenecks}

\subsection{Multi-energy systems}

The analysis of MES becomes increasingly important in energy system
modeling due to the ongoing electrification of the heating, transport,
and industry sectors, the shift towards decentralized energy production,
and the intermittent nature of renewable energy systems (RES) that
requires substantial investments in storage capacities and grid
extension. Although a clear definition of MES is absent, MES can be
understood as an extension of hybrid renewable energy systems (HRES),
where HRES usually denote systems combining two or more renewable
generation and storage technologies for reliable renewable electricity
supply, such as PV, Wind, and battery storage, while MES broaden this
scope to coupled electricity, heating, cooling, industry and transport
sectors through centralized and decentralized generation, storage and
conversion technologies, as illustrated in Figure
\ref{fig:multienergysystem}. This cross-sectoral assessment is
crucial owing to the ongoing electrification of the transport and
heating sectors and the intermittency challenge of renewable energy that
necessitates the provision of system flexibilities of a resilient,
cost-efficient, and decarbonized energy system.

While the definition of MES and integrated energy systems (IES) is often
used synonomously, MES refers to the technical modeling of sector
coupling (electricity, heating, transport, industry) and physical energy
flows between sectors. IES, on the other hand, encompass not only the
physical modeling but also regulatory and systemic operational
integration of energy sectors, infrastructures, and markets with the aim
of coordinated planning and operation of all RES and infrastructures.
MES builds thus the physical basis of IES which is the overarching
systemic concept for planning of including grid expansion, market
integration, and energy policies.

Coupling electricity to heat, gas, hydrogen, transport and industrial
processes in optimization models increases both the number of variables
and the number of non-linear relationships, and creates computationally
expensive problems for which surrogates are increasingly used.
\cite{Zhou20181,Zheng20231907,Ren2024,Jalving2023,Dong2023,Yin2024,Cao2023,Wang2020202933}.

\subsection{Surrogate models}

\begin{figure*}[!t]
\centering
\includegraphics[width=\textwidth]{./Surrogate_classification.png}
\caption{Conceptual positioning and classification of surrogate models in energy-system optimization. The figure illustrates how surrogate models mainly act as predictive approximations of expensive energy-system models, but acquire a prescriptive role when embedded in optimization workflows. The lower panel summarizes the four classification axes used in this review: model family, surrogate target, optimizer role and trust mechanism.}
\label{fig:surrogate-classification}
\end{figure*}

Surrogate models, also referred to as metamodels or emulators, are
computationally cheap approximations of an expensive original model.
They are constructed from input--output evaluations of the original
model and are then used to replace, approximate or support selected
parts of an analysis or optimization workflow
\cite{Bhosekar2018,Westermann2019,FernandezGodino2023}. In energy system
optimization, the expensive original model may be a detailed building
simulation, MES optimization, an optimal-power-flow formulation, a
mixed-integer dispatch model, or another high-fidelity numerical model
\cite{Westermann2019,Hirvonen2017323,Manco2024,Klemm2021,Ruan2021221}.
The purpose of the surrogate is therefore not to replace the physical or
optimization model, but to reduce the number, duration or complexity of
expensive model evaluations during design, operation, uncertainty
analysis or multi-objective search
\cite{Bhosekar2018,DiazManriquez2016,Hirvonen2017323}.

In the basic workflow, a set of representative system configurations is
sampled at first. Second, the original model is evaluated for these
configurations, producing training data. Third, a statistical or
machine-learning model is fitted to approximate the mapping from inputs
to outputs based on the generated training data. Once validated, the
surrogate can be evaluated much faster than the original model and can
therefore be used for repeated evaluations inside optimization,
sensitivity or uncertainty analysis
\cite{Bhosekar2018,Westermann2019,FernandezGodino2023}. In building
applications, for example, detailed simulation models, e.g. TRNSYS or
EnergyPlus, are used to generate training samples, while surrogate
models are then used to accelerate optimization
\cite{Hirvonen2017323,Westermann2019,Elwy2024,Zhang2026_building}. In
the solar-community study by Hirvonen et al., a neural network-based
surrogate model is updated iteratively during multi-objective
optimization where candidate solutions are predicted by the surrogate
and re-evaluated in the full TRNSYS model and subsequently added to the
training set \cite{Hirvonen2017323}. Surrogate modeling can either be an
offline approximation step or a part of the optimization loop itself
(in-loop).

\subsection{Optimization problem classes and computational bottlenecks}
\label{sec:optimization-bottlenecks}

Surrogate models become relevant where energy-system optimization
problems are repeatedly solved, computationally expensive, or physically
too detailed to be evaluated multiple times inside a search loop at high
spatio-temporal resolution. The relevant problem classes include
short-term dispatch and unit commitment, optimal power flow, capacity
and generation expansion, district-heating and MES design, stochastic or
robust planning considering uncertainty, and highly constrained
multi-objective optimization
\cite{Wu2022231,Chen20244723,Chen2022,Muhlpfordt20194806,Deng2020831,Khaloie2025,Rodgers2018951,Han2018622,Hirvonen2017323,Zhou20181}.

The computational burden arises from several sources: Dispatch and
models become difficult because of intertemporal storage constraints,
binary on/off decisions, and start-up costs of dispatchable sources.
Moreover, network-constrained problems such as AC OPF add nonlinear
voltage, reactive-power and security constraints require further
computational resources. Capacity-expansion, microgrid-sizing and
MES-design problems are expensive because each candidate solution
usually requires multiple dispatch evaluations. Sector coupling in MES
increases problem size through electricity-heat-gas-hydrogen conversion
technologies, storage State-of-charge (SoC) states and nonlinear
relationships
\cite{Safta20172535,Pang2023,Su2022315,Chen2023657,Koirala20242284,Zheng20231907,Ren2024,Jalving2023,Dong2023,Yin2024}.

Additionally, if uncertainty is considered, stochastic
chance-constrained formulations add scenarios, recourse decisions or
probabilistic constraints, while MOO problem create outer loops in which
multiple candidate solutions must be evaluated to approximate Pareto
fronts
\cite{Hu2020,Hu2021671,Volodina2022,Liang20246508,Thrampoulidis2021,Hussien20211883,Schulte2020,Han2018622}.
